\newpage
\section{Acta Número 10}
\label{sec:acta10}

\noindent \textbf{Fecha de la sesión:}\par
Lunes, 06 de Abril de 2026

\vspace*{0.5cm}
\noindent \textbf{Datos de la Sesión:}
\vspace{-0.2cm}
\begin{itemize}[noitemsep]
    \item \textbf{Modalidad:} Virtual - Google Meet.
    \item \textbf{Hora de Inicio:} 18:00
    \item \textbf{Hora de Fin:} 20:00
\end{itemize}

\vspace*{0.5cm}
\noindent \textbf{Orden del Día:}
\vspace{-0.2cm}
\begin{itemize}[noitemsep]
    \item Asistencia.
    \item Kick-off de la fase de Aseguramiento de Calidad (QA) del Incremento.
    \item Asignación estratégica de responsables QA por células de desarrollo.
    \item Definición de Definition of Done (DoD) para QA: Test Cases y Bug Reports.
    \item Daily Sync extendido: Impedimentos en Backend, Frontend y Base de Datos.
    \item Gestión del Sprint Backlog: Identificación de bloqueos y Carry Over.
    \item Refinamiento técnico sobre lineamientos del repositorio de documentación.
    \item Firmas y Aprobación de los Presentes.
\end{itemize}

\vspace*{0.5cm}
\noindent \textbf{Socios Fundadores Presentes:}
\vspace{-0.2cm}
\begin{enumerate}[noitemsep]
    \item Pardo Pozo Juan Diego (Jefe de Proyecto).
    \item Arnez Jaimes Mauricio.
    \item Quispe Flores Camila Belen.
    \item Ariñez Bautista Jim Gabriel.
    \item Medina Rodríguez Mateo Jhoustin.
    \item Molina Maldonado Tomas Manuel.
    \item Zeballos Crespo Jonathan.
    \item Gutierrez Guzmán Angheline.
\end{enumerate}

\newpage
\subsection{Desarrollo del Acta}

\noindent \textbf{Control de Asistencia del Team}\par
Se procedió a la toma de asistencia a la hora acordada, corroborando la presencia en sesión de los 8 socios fundadores de la grupo-empresa Byte Busters S.R.L. Al contar con la totalidad de los miembros, se estableció el quórum necesario para iniciar la toma de decisiones y el trabajo de planificación.

\vspace*{0.5cm}
\noindent \textbf{Kick-off de la fase de Aseguramiento de Calidad (QA) del Incremento}\par
Se oficializó el inicio de la etapa de validación de calidad. El objetivo principal de esta fase es someter el Incremento de producto desarrollado a rigorosas pruebas para asegurar su estabilidad antes de la potencial liberación.

\vspace*{0.5cm}
\noindent \textbf{Asignación estratégica de responsables QA por células de desarrollo}\par
Para garantizar la objetividad en la validación, se establecieron roles de revisión cruzada. Se designaron pares de QA responsables de auditar el trabajo de células homólogas, mitigando el sesgo de autoría en el código.

\vspace*{0.5cm}
\noindent \textbf{Definición de Definition of Done (DoD) para QA: Test Cases y Bug Reports}\par
Se estandarizó el marco de trabajo para la documentación de pruebas. Como parte de la actualización de nuestra Definition of Done (DoD), se exigirá a cada integrante el diseño y ejecución de un mínimo de 10 Test Cases. Asimismo, toda anomalía detectada deberá registrarse formalmente en el documento mediante Bug Reports.

\vspace*{0.5cm}
\noindent \textbf{Daily Sync extendido: Impedimentos en Backend, Frontend y Base de Datos}\par
Se abrió el espacio técnico para revisar el estado de los repositorios y el código base. Al verificar las capas de Frontend, Backend y Database, el equipo de desarrollo no reportó deudas técnicas urgentes ni bloqueos (Blockers) activos, permitiendo concluir este punto de forma ágil.

\vspace*{0.5cm}
\noindent \textbf{Gestión del Sprint Backlog: Identificación de bloqueos y Carry Over}\par
Tras la evaluación de viabilidad de las historias activas, se determinó que los ítems IDEN-007 (Herramienta de Moderación y Auditoría - Admin Tool) e IDEN-008 (Optimización Técnica, SEO y Accesibilidad del Hero) superan la capacidad restante del equipo debido a dependencias técnicas. En consenso, se decidió catalogarlas como \textit{Carry Over} y reasignarlas al Sprint Backlog del Sprint 2.

\vspace*{0.5cm}
\noindent \textbf{Refinamiento técnico sobre lineamientos del repositorio de documentación}\par
Se llevó a cabo una alineación para resolver dudas operativas sobre el entorno de documentación del proyecto. Se estandarizó el uso de archivos de configuración, ejecución de macros, paquetes, estructuración del layout y aplicación de la paleta de colores corporativa para asegurar la homogeneidad técnica en los entregables.

\newpage
\noindent \textbf{Aprobación Oficial}\par
Habiendo estructurado el plan de iteración y establecido los estándares operativos para el control de calidad, se da por aprobada la sesión, comprometiendo al equipo con los objetivos trazados.

\vspace*{0.5cm}
\noindent \textbf{Firmas y Aprobación de los Presentes}
\begin{center}
    \noindent
    \setlength{\tabcolsep}{0pt}
    \begin{tabular}{ccc}
        \espacioFirma{assets/img/firmas/mauricio.jpg}{Mauricio Jaimes Arnez}{13751221} & 
        \espacioFirma{assets/img/firmas/camila.jpg}{Camila Belen Quispe Flores}{13097244} &
        \espacioFirma{assets/img/firmas/jim.jpg}{Jim Gabriel Ariñez Bautista}{9517642} \\[1.0cm]
        \espacioFirma{assets/img/firmas/mateo.jpg}{Mateo Medina Rodríguez}{9453809} &
        \espacioFirma{assets/img/firmas/tomas.jpg}{Tomas Molina Maldonado}{8051701} &
        \espacioFirma{assets/img/firmas/jonathan.jpg}{Jonathan Zeballos Crespo}{13067494} \\[1.0cm]
    \end{tabular}
    \vspace{0.2cm} 
    \begin{tabular}{cc}
        \espacioFirma{assets/img/firmas/diego.jpg}{Juan Diego Pardo Pozo}{12747783} & 
        \espacioFirma{assets/img/firmas/angheline.jpg}{Angheline Gutierrez Guzmán}{13751815}
    \end{tabular}
\end{center}

\vspace{0.5cm}
\noindent \textbf{Fecha de última revisión:} Lunes, 06 de Abril de 2026